%==============================================================================
% CONCLUZII - nu se numeroteaza, apare ca element separat dupa Capitolul 4 si
% inainte de Bibliografie. Conform Regulamentului ETTI, Cap. 3, pct. j.
%==============================================================================

\chapter*{Concluzii}
\addcontentsline{toc}{chapter}{Concluzii}

Lucrarea de față a propus și a realizat ETTIway, o platformă web pentru
navigarea interactivă într-un campus universitar, plecând de la o problemă
reală, lipsa unei soluții digitale unitare pentru orientarea în campusul ETTI.
Totodată, platforma integrează identificarea clădirilor, căutarea sălilor și calculul
traseelor între puncte relevante. Această secțiune sintetizează rezultatele
obținute, formulează opinia personală asupra acestora și schițează direcții
naturale de dezvoltare ulterioară.

\section*{Sinteza rezultatelor}
\addcontentsline{toc}{section}{Sinteza rezultatelor}

Din punct de vedere tehnic, lucrarea a condus la o platformă funcțională,
structurată pe trei niveluri logice clar separate: o interfață web construită cu
JavaScript și Leaflet, un backend realizat cu Spring Boot și o bază de date
PostgreSQL pentru persistarea informațiilor. Au fost implementate
funcționalitățile centrale, cum ar fi autentificarea cu
diferențierea rolurilor de utilizator și administrator, editarea hărții direct
din interfață cu persistarea modificărilor în baza de date si căutarea contextuală
a sălilor și clădirilor. Printre care se numără și afișarea detaliilor asociate locațiilor și calculul
vizual al traseelor optime între puncte, bazat pe algoritmul lui Dijkstra,
combinat cu poziționarea în timp real a utilizatorului prin GPS. Dincolo de
navigarea exterioară, platforma include și o componentă de rutare în interiorul
clădirilor, în care administratorul poate desena planuri de etaj cu camere,
coridoare și noduri de rutare. Ulterior, utilizatorul poate calcula traseul între
două puncte de pe același etaj. Soluția folosește același algoritm Dijkstra
aplicat unui graf construit din planurile salvate ca documente GeoJSON, ceea
ce a permis reutilizarea codului existent fără modificări semnificative.

În planul infrastructurii, soluția adoptată pentru punerea în funcțiune se
distinge prin alegerea de a folosi o rețea privată virtuală construită peste
WireGuard, în locul expunerii clasice a serverului pe internetul public. Această
decizie de proiectare elimină complet suprafața de atac externă a sistemului. Ulterior vom putea 
expune platforma pentru publicul larg mult mai ușor.


La nivel metodologic, lucrarea a parcurs ciclul complet de dezvoltare a unei
aplicații reale: analiza cerințelor, proiectarea arhitecturii, implementarea
incrementală, integrarea componentelor, punerea în funcțiune într-un mediu
accesibil de pe mai multe dispozitive și testarea sistematică în condiții reale
de utilizare. Această experiență, mai mult decât rezultatul tehnic în sine,
constituie o componentă importantă a contribuției personale, expunându-mă la
deciziile concrete care apar la fiecare nivel al unui sistem software complet.


\newpage

\section*{Opinia personală asupra rezultatelor}
\addcontentsline{toc}{section}{Opinia personală asupra rezultatelor}

Consider că obiectivele declarate ale lucrării au fost atinse în măsură
satisfăcătoare. Platforma este funcțională în mediul real, oferă o experiență de
navigare fluidă în spațiile exterioare ale campusului și demonstrează că o
soluție de orientare digitală poate fi construită folosind tehnologii deschise,
fără dependențe comerciale sau costuri de licențiere. Combinarea backend-ului
Java cu o interfață web bazată pe JavaScript și Leaflet s-a dovedit potrivită
pentru scopul propus, oferind atât performanța necesară, cât și o experiență de
utilizator comparabilă cu cea a aplicațiilor de navigare cunoscute.

Cea mai valoroasă lecție a procesului de dezvoltare a fost legată de
importanța deciziilor de infrastructură. Inițial, atenția mea s-a concentrat pe
funcționalitățile vizibile ale aplicației, harta, căutarea, rutarea, însă în
practică s-a dovedit că punerea în funcțiune sigură și corectă a platformei
ridică probleme la fel de interesante și solicită soluții la fel de gândite. În
particular, înțelegerea legăturii dintre contextul securizat HTTPS și accesul la
poziția GPS a fost una dintre observațiile concrete ale proiectului, care a
modelat decisiv alegerea soluției de punere în funcțiune.

În același timp, recunosc că platforma, în starea sa actuală, are limite clare.
Acoperirea datelor este parțială, iar precizia poziționării în interiorul
clădirilor este afectată de constrângerile tehnologice ale GPS-ului, descrise
în Capitolul 1 și confirmate experimental în Capitolul 4. Aceste limitări nu
sunt eșecuri ale soluției, ci delimitări naturale ale domeniului său actual de
aplicabilitate, iar identificarea lor explicită deschide direcții firești de
dezvoltare ulterioară.

\section*{Direcții viitoare de dezvoltare}
\addcontentsline{toc}{section}{Direcții viitoare de dezvoltare}

Pe baza experienței acumulate și a limitărilor identificate, mai multe direcții
de dezvoltare apar ca firești și valoroase.

Prima direcție vizează extinderea funcționalității de navigare indoor către
deplasări între etaje diferite. Versiunea actuală tratează rutarea pe un
singur etaj la un moment dat, iar selectarea unui punct de plecare și a unei
destinații aflate pe etaje diferite este semnalată ca neimplementată.
Realizarea acesteia ar presupune adăugarea unor muchii speciale între nodurile
de tip scară de pe etaje succesive, precum și a unei logici de tranziție
vizuală între planurile etajelor pe măsură ce traseul calculat le traversează.
O extindere conexă privește integrarea fluxului exterior–interior: în
versiunea actuală, utilizatorul comută explicit între harta campusului și
planul unei clădiri, dar un traseu unificat ar putea porni de la o poziție
GPS exterioară, ar putea trece prin intrarea principală a clădirii și ar
putea continua pe coridoarele indoor până la sala destinație. Pentru
poziționarea efectivă a utilizatorului în interiorul clădirilor, unde GPS-ul
devine nefiabil, ar putea fi explorate soluții complementare bazate pe balize
Bluetooth de joasă energie sau pe semnale Wi-Fi.

A doua direcție privește îmbogățirea acoperirii de date. Adăugarea de noi
clădiri, săli, alei și puncte de interes este deja posibilă din interfața de
administrare, fără modificări în cod, deci această direcție nu necesită
dezvoltare tehnică, ci muncă de întreținere a datelor. Pe termen mediu, ar fi
utilă explorarea unor mecanisme de import semiautomat din surse externe, precum
OpenStreetMap, pentru a accelera această muncă.

A treia direcție ține de funcționalități avansate de rutare. Algoritmul lui
Dijkstra implementat oferă traseul de cost minim, considerând doar distanța, dar
ar putea fi extins pentru a ține cont și de alte criterii: accesibilitate pentru
persoane cu dizabilități, evitarea zonelor aglomerate, preferința pentru
traseele acoperite în caz de vreme nefavorabilă. Asemenea extinderi presupun
îmbogățirea grafului cu informații suplimentare la nivelul muchiilor și o
interfață prin care utilizatorul să își exprime preferințele.

Una dintre direcțiile care depășește perimetrul strict tehnic este testarea pe
scară mai largă, cu utilizatori reali din comunitatea universitară.
O astfel de evaluare ar oferi informații pe care testarea individuală nu le
poate înlocui, privind ergonomia interfeței, acoperirea reală a nevoilor de
navigare și acceptarea soluției ca instrument util în viața de zi cu zi în
campus.

\section*{Reflecție finală}
\addcontentsline{toc}{section}{Reflecție finală}
 
Dincolo de rezultatul tehnic, lucrarea a reprezentat o experiență în
care am avut ocazia să parcurg, în mod integrat, etape pe care, în mod normal,
le abordăm separat pe parcursul anilor de studiu: proiectare software,
programare web, baze de date, securitate, administrare de sisteme. Această
integrare, dificilă la început, mi-a oferit o înțelegere mai matură a modului în
care o soluție tehnică reală se construiește prin succesiunea unor decizii la
mai multe niveluri, fiecare cu propriile sale considerații și compromisuri.
 
Pe parcursul proiectului am întâlnit mai multe situații care m-au făcut să
înțeleg mai bine echilibrul dintre alegerile teoretice ideale și constrângerile
practice. Decizia de a stoca graful campusului sub formă de JSON serializat,
și nu într-o schemă relațională cu mai multe tabele, a fost luată după o
analiză a complexității necesare versus flexibilitatea câștigată. În
cele din urmă, această decizie s-a dovedit corectă, pentru că extinderea către
componenta indoor s-a făcut natural, prin adăugarea unei a doua coloane JSON,
fără modificări de schemă. La fel, alegerea de a folosi JavaScript vanilla
în loc de un framework modern precum React a impus o disciplină mai mare în
organizarea codului, dar a păstrat dependențele externe minime și mi-a permis
să înțeleg în profunzime fiecare mecanism al aplicației, fără să mă bazez pe
capacitățile unui framework.
 
Una dintre cele mai surprinzătoare lecții a fost legată de aspectele
non-funcționale ale proiectului. Inițial, atenția mea s-a concentrat pe
funcționalitățile vizibile: cum desenez harta, cum implementez căutarea, cum
calculez traseul. Pe parcurs am descoperit că la fel de importante sunt
deciziile despre cum punem aplicația în funcțiune, cum o protejăm de accese
neautorizate, cum o facem accesibilă de pe dispozitive mobile reale.
Descoperirea că accesul la poziția GPS din browser cere obligatoriu un context
securizat HTTPS a fost un moment în care am înțeles că proiectarea unei
aplicații web moderne presupune o cunoaștere a infrastructurii la fel de
solidă ca a algoritmilor de bază. Această observație a modelat decisiv
alegerea de a folosi Tailscale, care a rezolvat simultan problema accesului
HTTPS automat și pe cea a expunerii sigure a serverului.
 
Lucrând la indoor, am realizat că soluțiile pentru spațiul interior diferă
fundamental de cele pentru exterior, în pofida similarității aparente.
Coordonatele pixel ale unui plan desenat manual nu pot fi tratate cu aceleași
funcții ca latitudinile și longitudinile reale; toleranța de \textit{snapping}
folosită pentru a conecta puncte aproximate la rețeaua de coridoare este o
soluție pragmatică, iar valoarea de 25 pixeli a fost stabilită experimental,
încercând mai multe valori până când rezultatul a devenit predictibil. Astfel
de detalii, care nu apar în cărțile despre algoritmi, sunt cele care fac
diferența între o demonstrație academică și o aplicație utilizabilă.
 
Consider că tocmai acest tip de experiență, în care am parcurs
toate straturile, de la algoritm la deployment, de la backend la mobile, de
la teoretic la practic, este cel mai prețios beneficiu pe care l-am obținut
prin elaborarea acestei lucrări. ETTIway nu este doar un produs software, ci
și o oportunitate de a verifica cunoștințele acumulate într-un proiect care
îmbină multe domenii și care, prin punerea sa în funcțiune în campus, devine
o soluție reală, nu doar o exercitare academică.
 